\lecture{43}{FedAvg}

\sourcevideo
    {https://www.youtube.com/playlist?list=PLOzRYVm0a65fhKDS8cYIC7YCuVGJ4FOJx}
    {Week 11: Lecture 43: FedAvg Algorithm}

FedAvg реализует основной цикл федеративного обучения: сервер
рассылает текущую модель, выбранные клиенты выполняют несколько
локальных шагов, после чего сервер усредняет полученные модели. В
отличие от обычного распределённого градиентного спуска, между двумя
коммуникационными раундами здесь выполняется несколько локальных
обновлений.

\section{Локальная и глобальная задачи}

Пусть имеется \(K\) клиентов. Клиент \(k\) хранит набор
\[
    \mathcal D_k
    =
    \{\zeta_{k,1},\ldots,\zeta_{k,n_k}\},
    \qquad
    n_k=\lvert\mathcal D_k\rvert,
\]
и локальную функцию
\[
    F_k(x)
    =
    \frac1{n_k}
    \sum_{j=1}^{n_k}
    \ell(x;\zeta_{k,j}),
    \qquad
    x\in\mathbb R^d.
\]
Если
\[
    n=\sum_{k=1}^{K}n_k,
    \qquad
    p_k=\frac{n_k}{n},
\]
то глобальная функция равна
\[
    F(x)
    =
    \sum_{k=1}^{K}p_kF_k(x),
    \qquad
    \sum_{k=1}^{K}p_k=1,
\]
а задача обучения имеет вид
\[
    \min_{x\in\mathbb R^d}F(x).
\]
Вес \(p_k\) учитывает объём данных клиента.

Алгоритм задаётся числом коммуникационных раундов \(T\), долей
участвующих клиентов \(C\in(0,1]\), размером локального мини-пакета
\(B\), числом локальных эпох \(E\) и шагом \(\eta>0\). В одном
раунде сервер выбирает
\[
    m=\max\{1,\lfloor CK\rfloor\}
\]
клиентов. Если клиент \(k\) за одну эпоху обрабатывает весь набор
\(\mathcal D_k\), то число его локальных шагов равно
\[
    \tau_k
    =
    E\left\lceil\frac{n_k}{B}\right\rceil.
\]
При одинаковых \(E\) и \(B\) клиенты с различными объёмами данных
могут выполнить разное число обновлений.

\section{Один раунд FedAvg}

На раунде \(t\) сервер выбирает множество
\[
    S_t\subseteq\{1,\ldots,K\},
    \qquad
    \lvert S_t\rvert=m,
\]
и отправляет каждому \(k\in S_t\) модель \(x^t\). Клиент
инициализирует
\[
    x_k^{t,0}=x^t
\]
и для \(s=0,\ldots,\tau_k-1\) выбирает мини-пакет
\(\mathcal B_k^{t,s}\subseteq\mathcal D_k\), вычисляет
\[
    g_k^{t,s}
    =
    \frac1{\lvert\mathcal B_k^{t,s}\rvert}
    \sum_{\zeta\in\mathcal B_k^{t,s}}
    \nabla_x\ell(x_k^{t,s};\zeta)
\]
и выполняет шаг
\[
    x_k^{t,s+1}
    =
    x_k^{t,s}-\eta g_k^{t,s}.
\]
После \(\tau_k\) шагов сервер получает модель
\(x_k^{t,\tau_k}\) или эквивалентное приращение
\[
    \Delta_k^t
    =
    x_k^{t,\tau_k}-x^t.
\]

Положим
\[
    n_{S_t}=\sum_{j\in S_t}n_j.
\]
Федеративное усреднение задаётся формулой
\[
    x^{t+1}
    =
    \sum_{k\in S_t}
    \frac{n_k}{n_{S_t}}x_k^{t,\tau_k}.
\]
Через приращения то же обновление записывается как
\[
    x^{t+1}
    =
    x^t
    +
    \sum_{k\in S_t}
    \frac{n_k}{n_{S_t}}\Delta_k^t.
\]
Если участвуют все клиенты, то \(n_{S_t}=n\), и серверные веса
совпадают с весами глобальной функции.

\section{Связь с градиентным спуском и локальный дрейф}

Предположим, что участвуют все клиенты и каждый выполняет один точный
градиентный шаг. Тогда
\[
    x_k^{t,1}
    =
    x^t-\eta\nabla F_k(x^t),
\]
и после усреднения
\[
\begin{aligned}
    x^{t+1}
    &=
    \sum_{k=1}^{K}p_k
    \bigl(x^t-\eta\nabla F_k(x^t)\bigr)
    \\
    &=
    x^t-\eta\nabla F(x^t).
\end{aligned}
\]
В этом частном случае FedAvg точно воспроизводит централизованный
градиентный шаг.

При \(\tau_k>1\) последующие градиенты вычисляются уже в различных
точках
\[
    x_k^{t,0},x_k^{t,1},\ldots,x_k^{t,\tau_k-1},
\]
поэтому серверное усреднение в общем случае не совпадает с одним
шагом по \(F\). Если локальные функции различаются, направления
\(\nabla F_i\) и \(\nabla F_j\) также различаются, и локальные модели
дрейфуют друг от друга. Увеличение числа локальных шагов сокращает
число коммуникаций, но усиливает этот дрейф; при неоднородных данных
он может заметно ухудшить направление глобального обновления.

\section{Эксперимент с MNIST}

Рассмотрим набор MNIST из \(60000\) изображений,
распределённых между \(100\) клиентами по \(600\) объектов. При
однородном разбиении изображения случайно перемешиваются, поэтому
локальные наборы приблизительно воспроизводят общее распределение
десяти классов.

Для неоднородного разбиения изображения сначала сортируются по меткам,
затем делятся на \(200\) блоков по \(300\) объектов, и каждый клиент
получает два блока. Поэтому его локальный набор обычно содержит только
две цифры. В экспериментах с полносвязной и свёрточной сетями такое
разбиение требовало большего числа коммуникационных раундов, а часть
настроек не достигала заданной точности. Причина состоит в том, что
локальный градиент строится по малому числу классов и может существенно
отклоняться от глобального направления.

\section{Параметры и стоимость}

Увеличение доли участия \(C\) даёт серверу более представительную
выборку клиентов и обычно уменьшает число требуемых раундов. При этом
число передаваемых моделей в одном раунде растёт примерно как
\(m\approx CK\), поэтому коммуникационная стоимость раунда
увеличивается.

Увеличение числа эпох \(E\) или уменьшение размера мини-пакета \(B\)
повышает число локальных шагов
\[
    \tau_k
    =
    E\left\lceil\frac{n_k}{B}\right\rceil.
\]
Это может сократить число обращений к серверу, но увеличивает локальную
вычислительную работу и потенциальный дрейф моделей. В таблицах
лекции обозначение \(B=\infty\) соответствует вычислению градиента по
всему локальному набору.

Следовательно, число коммуникационных раундов само по себе не
определяет эффективность FedAvg. Необходимо учитывать размер модели,
число участвующих клиентов, объём локальных вычислений, доступность
устройств и степень неоднородности данных.

\section{Ограничения базовой схемы}

Базовый FedAvg предполагает, что выбранные клиенты завершают локальное
обучение и возвращают корректные обновления. Он непосредственно не
устраняет задержки медленных клиентов, разрывы связи, злонамеренные
обновления, утечку информации из локальных моделей и сильный дрейф
при неоднородных данных. Для этих случаев требуются дополнительные
механизмы отбора клиентов, робастной и защищённой агрегации,
приватности и коррекции локального обучения.

Главный компромисс FedAvg состоит в том, что несколько сравнительно
дешёвых локальных шагов выполняются между двумя дорогими обменами с
сервером. Поэтому параметры \(C\), \(E\), \(B\) и \(\eta\) должны
выбираться совместно, а не независимо.
